\section{Heat-flow leakage and the pin bootstrap}\label{sec:leakage}

We now use the geometric estimates to bound the return fraction in the
simultaneous pin test.  The main task is to prove that the hitting
function loses nearly $m$ units of squared norm during each of two short
portions of a period, where $m$ is the number of new pins.  Local diffusion
spreads the mass on those pins, and the flow estimate controls its pairing
with the rest of the hitting function.  Averaging the accumulated heat
dissipation will absorb this pairing.
The Nash method for Markov semigroups is developed in
\cite[Section~3]{DiaconisSaloffCoste1996}.  The proof below accounts
explicitly for the changing hole sets, killing, instantaneous
contractions, and the restriction of the Nash inequality to regular
times.

Throughout the section, $d\ge3$, $\ell=2S$, and $\Lambda$ is a finite cube.
We use the parameters of Section~\ref{sec:sparse}, namely
\[
 \alpha=\frac1{100},\qquad \nu=\frac1{100d},\qquad
 a=\frac{11}{10},\qquad s=p^{-a},\qquad R=\lceil C_*s\rceil,
 \qquad \beta\ge5s.
\]
The spatial growth constant $L$, the regularity radius $M_*$, and the
constant $C_*$ will be fixed independently of the final choice of $p$.
The test $T$ consists of $m\ge1$ distinct slot points at one time, is
disjoint from $E\supseteq D_0$, and satisfies
\begin{equation}\label{eq:leakage-growth}
 \#\{i\in T:x_i\in B_h(z)\}\le L(1+h)^\alpha
 \quad(z\in\Z^d,\ h\ge0).
\end{equation}
Here $x_i$ denotes the site of the slot point $i$.
Write $P=E\cup T$, and sample the exposed paths under $\mathbb P_P$,
as in Proposition~\ref{prop:exposure}.  Expectations of quantities
determined by these paths use their actual marginal, including the factor $g_P$ in
\eqref{eq:exposure-outer-law}.  At a fixed exposure, the propagators use
the unconditioned fresh-turn walk.  As in
Proposition~\ref{prop:pin-test}, $r$ is the
expected fraction of points of $T$ whose next pin along the actual
cycle again belongs to $T$.

Recall the one-period, column-forward kernel $U$ of the fresh-turn walk,
killed on visiting $D_0$.  At the cut containing $T$, the function $f$
of Lemma~\ref{lem:return-energy} is $1$ on $T$ and, off $T$, is the
probability of a future visit to $T$ before killing, namely the minimal
nonnegative hitting solution.
In particular,
\begin{equation}\label{eq:leakage-hitting-input}
 0\le f\le1,\qquad f|_{\partial\Lambda}=0,\qquad
 U^*f=f\ \hbox{off }T,\qquad
 m(1-r)\ge\mathbb E_P\langle f,(I-U)f\rangle.
\end{equation}
The last inequality incorporates the conditional first-moment identity
of Proposition~\ref{prop:first-moment} and the relaxation of the other
pins of $E$.  A class of the relaxed walk that never reaches $T$ has
$f=0$; certain absorption is unnecessary.  In particular $f$ is a
function of the exposed paths, and we will not require it to be
independent of any geometric event.
All the propagators below are measurable functions of finite-jump
data.  The hitting function is measurable as the increasing limit of
finite-horizon hitting probabilities, so the expectations used below
are well defined.

For grid pin sets satisfying the larger-pin induction hypothesis
$\mathsf I_p(P)$ of \eqref{eq:induction-assertion}, our aim is
\begin{equation}\label{eq:leakage-target}
 \mathbb E_P\langle f,(I-U)f\rangle\ge m[1-\varepsilon(p)],
 \qquad \varepsilon(p)=o(p).
\end{equation}
Together with \eqref{eq:leakage-hitting-input}, this gives
$r\le\varepsilon(p)$.  The simultaneous pin test then recovers
$\mathbb P_E(T\subseteq\mathcal B_E)\le p^m$ once
$2\varepsilon(p)\le p$, closing the induction step.
To obtain \eqref{eq:leakage-target}, split
$f=\eta+g$, where $\eta=\1_T$ and $g=f-\eta$.
Then $0\le g\le1$ and $\langle\eta,g\rangle=0$.
We first relate the quadratic form to norm losses at both ends of the
period.  The Nash estimate will make the propagated $\eta$ small in
squared norm, while the flow estimate will control its pairing with the
propagated $g$ through the square root of that function's energy at the
same time.

\subsection{Ordered propagation and the two ends of a period}

We first fix one exposed environment.  Between its finitely many
instantaneous changes, the surviving interior holes carry the symmetric
rate-$1/2$ heat equation.  A boundary hole is assigned value zero.  Thus
the generator on an interior hole $k$ is
\begin{equation}\label{eq:leakage-generator}
 (L_t v)_k=\frac12\sum_{j:\{j,k\}\text{ is an available slot edge}}
                   (\widetilde v_j-v_k),
\end{equation}
where $\widetilde v$ equals $v$ on interior holes and zero on boundary
holes.  The diagonal includes the escape rates to the boundary.  This
is continuous killing throughout the heat interval.

At an instantaneous change, let $F$ be the forced bijection or the
average of the allowed within-site completion bijections, from the
hole set just before the change to that just after it.  If $R_-$ and
$R_+$ are the projections deleting killed coordinates on those two
sets, respectively, the killed factor is
\begin{equation}\label{eq:ordered-killing}
             R_+FR_-,\qquad
             (R_+FR_-)^*=R_-F^*R_+.
\end{equation}
The projections remain in this order.  For heat intervals, the
corresponding assertion follows either from the symmetric Dirichlet
generator or by transposing the sampled surviving paths.

By the subpermutation argument in \eqref{eq:hole-contractions}, these
matrices, their products, and their adjoints are nonnegative, have row
and column sums at most one, and contract the counting $\ell^2$ norm.
We will also use the two sum bounds separately: the row bound gives
$0\le Av\le1$ when $0\le v\le1$, while the column bound gives
$\|Av\|_1\le\|v\|_1$ for $v\ge0$.

Put the time of $T$ at the beginning and end of a period.  Let $C_u$ be
the initial propagation of duration $u$, and let $C_u^{\rm rev}$ be
the adjoint of the ending propagation of duration $u$.
Forward propagation includes the changes after its initial cut through
its final cut, under the post-completion convention at time zero.
Reverse propagation traverses precisely those factors in the opposite
order.  If the ending cut is a completion cut, its adjoint completion
is the first factor traversed in reverse time.  It is included exactly
once, including when it occurs immediately at elapsed time zero.
We set $C_0=C_0^{\rm rev}=I$ on the initial, already unkilled
coordinates and count any such initial reverse jump in every positive
duration that traverses it.  This convention allows a contraction jump
at $u=0+$.

\begin{lemma}[Leakage at both ends]\label{lem:two-ends}
For every fixed exposed environment, $s>0$, and period $\beta\ge4s$,
the preceding killed propagators satisfy, for every real vector $v$
on the holes at the distinguished cut with zero boundary values,
\begin{equation}\label{eq:two-ends}
 \langle v,(I-U)v\rangle\ge\frac12\left(
   \|v\|_2^2-\|C_{2s}v\|_2^2
  +\|v\|_2^2-\|C_{2s}^{\rm rev}v\|_2^2\right).
\end{equation}
Both $u\mapsto\|C_uv\|_2$ and
$u\mapsto\|C_u^{\rm rev}v\|_2$ are nonincreasing.
\end{lemma}

\begin{proof}
The nonoverlapping initial and ending portions factor $U=BMA$, with
$A=C_{2s}$, $B^*=C_{2s}^{\rm rev}$, and $M$ the intervening
propagation.  At $\beta=4s$, assign every shared-cut jump once according
to the preceding convention and take the intervening identity on that
cut.  All factors are contractions, so
\[
 \langle v,Uv\rangle
 =\langle B^*v,MAv\rangle
 \le\|B^*v\|_2\|Av\|_2
 \le\tfrac12\bigl(\|B^*v\|_2^2+\|Av\|_2^2\bigr).
\]
This proves \eqref{eq:two-ends}.  Longer initial portions are obtained
by left multiplication by further contraction factors.  For the ending
adjoints the same fact follows after reversing the factor order, using
\eqref{eq:ordered-killing}.  It proves both monotonicity assertions.
\end{proof}

The reverse construction has symmetric heat between changes, inverse
forced transports, averages of inverse completion bijections, and the
ordered killing projections.  Thus it has the same deterministic
analytic properties as forward propagation.  Its environment is the
original exposed environment read in reverse order; no invariance of
the environment law under time reversal is asserted or needed.

Thus \eqref{eq:leakage-target} will follow if each of these two
propagations loses at least $m[1-\varepsilon(p)]$ in expected squared
norm by time $2s$.  We prove both estimates through the same argument.

For the rest of the estimates, $C_u$ denotes either of the two
propagations, and $t(u)$ denotes its physical time slice modulo $\beta$.
On the available slot graph, extended by all slots outside $\Lambda$,
use the energy
\[
 \mathcal D_{t}(v)=\frac12\sum_{\{j,k\}\text{ available at }t}
                          (v_j-v_k)^2.
\]
Vectors propagated in the box are extended by zero on its boundary and
outside.  Since the boundary is the inner boundary, any edge crossing
out of the box has two zero values.

Write $\eta_u=C_u\eta$,
$g_u=C_ug$, and
$\mathcal L(u)=\|f\|_2^2-\|C_uf\|_2^2$.
For a jump traversed at elapsed time $v$, let
\[
 \Delta_g(v)=\|g_{v,\mathrm{before}}\|_2^2
                  -\|g_{v,\mathrm{after}}\|_2^2\ge0.
\]
Here ``before'' and ``after'' refer to the chosen direction of
propagation.  Denote the finite list of jumps traversed by duration
$u$ by $\mathcal J(u)$, including an initial reverse jump when present.
Between jumps, \eqref{eq:leakage-generator} gives
\[
 \frac{\mathrm d}{\mathrm du}\|g_u\|_2^2
     =-2\mathcal D_{t(u)}(g_u).
\]
To check the normalization, an interior--interior edge contributes
$-(g_j-g_k)^2$ to this derivative, and an interior--boundary edge
contributes $-g_j^2$.  These are exactly the contributions to
$-2\mathcal D$.  Summing continuous and instantaneous losses and then
expanding $\|\eta_u+g_u\|_2^2$ yields the exact identity
\begin{align}
 \mathcal L(u)
 &=m-\|\eta_u\|_2^2-2\langle\eta_u,g_u\rangle\notag\\
 &\quad+2\int_0^u\mathcal D_{t(v)}(g_v)\,\mathrm dv
       +\sum_{v\in\mathcal J(u)}\Delta_g(v).
       \label{eq:exact-loss}
\end{align}
The sum includes a completion at an endpoint whenever that completion
is a factor of the propagation.  Dropping it gives
\begin{equation}\label{eq:loss-lower}
 \mathcal L(u)\ge
 m-\|\eta_u\|_2^2-2\langle\eta_u,g_u\rangle
       +2\int_0^u\mathcal D_{t(v)}(g_v)\,\mathrm dv.
\end{equation}
Keeping this as an inequality is necessary even when a completion is
the only nontrivial instantaneous operation.

\subsection{Localization with actual and fresh marks}

For $i\in T$, let $\eta_{i,u}$ be the propagation from unit mass at $i$,
with the additional rule that a path is killed when its site leaves
$B_R(x_i)$.  Use this additional killing for continuous moves and on
both sides of instantaneous changes.  Positivity and the path
description give
\begin{equation}\label{eq:local-pieces}
 0\le\zeta_u:=\sum_{i\in T}\eta_{i,u}\le\eta_u\le1,
 \qquad \|\eta_{i,u}\|_1\le1.
\end{equation}

\begin{lemma}[Localization averaged over the pin law]\label{lem:localization}
Fix $d$ and $\ell$.  There is a constant $C_*=C_*(d,\ell)$ such that
for $s\ge1$, $R=\lceil C_*s\rceil$, $\beta\ge5s$, every finite
cube, every $E\supseteq D_0$ consisting of $D_0$ and finitely many
pins, and every disjoint single-time test $T$, one has
\begin{equation}\label{eq:localization}
 \mathbb E_P\|\eta_u-\zeta_u\|_1\le C m e^{-cs}
       \qquad(0\le u\le2s)
\end{equation}
for either direction of propagation.  The constants are independent of
the volume, the pin sets, $\beta$, and $s$.
\end{lemma}

\begin{proof}
First sample a complete actual exchange picture with law
$\mathbb P_P$ and expose its missing cycles.  Conditional on that
exposure, realize the unconditioned one-particle propagation of
Proposition~\ref{prop:exposure} using independent fresh rate-$1/2$
Poisson clocks on all slot edges, accepting a fresh mark only when its
two endpoints are holes.  The forced moves use marks of the actual
picture.  Completion randomness changes slots only within a site.
This construction is for the unconditioned fresh walk; its exposure
still has the marginal induced by the full pin law.

The difference in \eqref{eq:local-pieces}, summed over output slots, is
bounded by the sum of the probabilities that these individual walks
exit their respective balls before elapsed time $u$.  Exiting requires
at least $R$ nearest-neighbor spatial steps.  The first $R$ such steps
give a chronological list of marks, each from either the actual
picture or the fresh family.  For a fixed list of sites, slot edges,
and choices of source, the actual marks used are distinct: the time
interval has length at most $2s<\beta$, so its image on the time circle
does not use a mark twice.  This remains true if the interval crosses
time zero, and if it is traversed backwards.  Fresh marks in the list
are distinct as well.

Suppose a source choice uses $k$ actual marks and $R-k$ fresh marks.
By Lemma~\ref{lem:factorial}, the joint factorial measure of the
actual marks is bounded by $2^k$ times their base product intensity.
Independence of the fresh family then bounds the mixed measure by
\[
        1^k(1/2)^{R-k}\,\mathrm dt_1\cdots\mathrm dt_R
\]
on the prescribed edge coordinates.  All source choices are integrated
over the same ordered simplex
$0<t_1<\cdots<t_R<u$, of volume $u^R/R!$.
At each spatial step there are at most $2d\ell^2$ choices for the next
site and the slot edge, even if one ignores constraints from the
current slot.  Summing source choices and edge choices therefore bounds
the expected number of such lists, for one initial point, by
\begin{equation}\label{eq:mixed-factorial-count}
 (2d\ell^2)^R(1+1/2)^R\frac{u^R}{R!}
       \le\frac{(2Ks)^R}{R!},\qquad K=3d\ell^2.
\end{equation}
The existence of an exit path is bounded by this count; the counted
lists need not themselves be realizable walk histories.  In
particular, the restrictions imposed by the exposure can be ignored
when taking this upper bound.  No factorial bound conditional on a
fixed exposure has been used.

Since $R!\ge(R/e)^R$, choosing $C_*>4eK$ bounds the last expression
by $2^{-R}\le e^{-cs}$.  Sum over $i\in T$ to obtain
\eqref{eq:localization}.  Transposition merely reverses the
chronological lists and the forced transports.  The same mixed
factorial estimate applies under the original law $\mathbb P_P$,
which proves the reverse assertion too.
\end{proof}

\subsection{Diffusion during regular times}

We now use the downward-induction hypothesis only through the three
proved inputs from Sections~\ref{sec:sparse} and~\ref{sec:flows}.
Assume $\mathsf I_p(P)$, as defined in
\eqref{eq:induction-assertion}.  With the fixed choices of constants
there, Proposition~\ref{prop:regularity} gives
\begin{equation}\label{eq:leakage-delta}
 \mathbb P_P(t\text{ is not regular for }i)
       \le\delta:=C R^d p^{2d}+C e^{-cR^\xi}
       \quad(i\in T, t\in[0,\beta)),
\end{equation}
where $\xi>0$ and the constants are fixed for all $p$ below a
preliminary threshold.  At a regular slice,
Proposition~\ref{prop:local-nash} states
\begin{equation}\label{eq:leakage-nash-input}
 \|v\|_2^{2+4/d}\le C\mathcal D_t(v)\|v\|_1^{4/d}
\end{equation}
for nonnegative $v$ supported on surviving slots above $B_R(x_i)$,
with zero extension.  Finally, Proposition~\ref{prop:flow-energy}
applies to any subset of the anchors regular at that same time and to
any nonnegative injections of mass at most one per anchor supported
above its $R$-ball.  Their sum is the divergence of a flow $J$ with
\begin{equation}\label{eq:leakage-flow-input}
                   \sum_e J(e)^2\le C mR^\alpha.
\end{equation}
The flow estimate follows from Proposition~\ref{prop:flow-decay} and
the growth condition; its applicability to arbitrary injection points
and subunit masses is part of its statement.

\begin{lemma}[Decay after enough regular time]\label{lem:regular-time-decay}
Let $T$ satisfy \eqref{eq:leakage-growth}, and assume
$\mathsf I_p(P)$ for $P=E\cup T$, with $p$ below the preliminary
threshold in Proposition~\ref{prop:regularity}.  For the localized
propagators defined above, in either direction and for
$s\le u\le2s$,
\begin{align}
 \mathbb E_P\|\eta_{i,u}\|_2^2
       &\le C(s^{-d/2}+\delta)\quad(i\in T),
                    \label{eq:piece-decay}\\
 \mathbb E_P\|\eta_u\|_2^2
       &\le C m\bigl[e^{-cs}+R^\alpha(s^{-d/2}+\delta)\bigr].
                    \label{eq:all-pieces-decay}
\end{align}
\end{lemma}

\begin{proof}
Put $q_i(u)=\|\eta_{i,u}\|_2^2$.  All localized factors are
contractions, so $0\le q_i\le1$ and $q_i$ is nonincreasing,
including at jumps.  Between jumps the additional killing gives the
Dirichlet generator with zero values outside $B_R(x_i)$ as well as
on the physical boundary.  Consequently, at every regular smooth
time, \eqref{eq:leakage-nash-input} and
$\|\eta_{i,u}\|_1\le1$ imply
\begin{equation}\label{eq:nash-differential}
 q_i'(u)=-2\mathcal D_{t(u)}(\eta_{i,u})
                         \le-c q_i(u)^{1+2/d}.
\end{equation}
At other smooth times $q_i'\le0$.
As long as $q_i>0$, the function $q_i^{-2/d}$ thus increases at a
fixed positive rate during regular time and cannot decrease at other
times or jumps.  If $q_i$ reaches zero, it stays zero.  It follows in
either case that
\[
 \int_0^s\1_{\{t(v)\text{ regular for }i\}}\,\mathrm dv\ge s/2
 \quad\Longrightarrow\quad
 q_i(u)\le C s^{-d/2}\quad(s\le u\le2s).
\]
The expected amount of nonregular time in $[0,s]$ is at most
$s\delta$ by \eqref{eq:leakage-delta} and Fubini.  Markov's
inequality bounds the probability of its exceeding $s/2$ by
$2\delta$.  This proves \eqref{eq:piece-decay}.  For reverse
propagation, $t(v)$ simply runs in the opposite direction; the same
fixed-time estimate \eqref{eq:leakage-delta} gives the same integral
bound, without changing the probability law.

At any slot above $z$, only anchors with $x_i\in B_R(z)$ can
contribute to $\zeta_u$.  Their number is at most
$L(1+R)^\alpha\le C R^\alpha$ by
\eqref{eq:leakage-growth}.  Hence pointwise Cauchy--Schwarz and
summation give
\[
                \|\zeta_u\|_2^2\le C R^\alpha\sum_{i\in T}q_i(u).
\]
Since $0\le\zeta_u\le\eta_u\le1$,
\[
 \|\eta_u\|_2^2-\|\zeta_u\|_2^2
       =\sum_k(\eta_u(k)-\zeta_u(k))(\eta_u(k)+\zeta_u(k))
       \le2\|\eta_u-\zeta_u\|_1.
\]
Use Lemma~\ref{lem:localization} and \eqref{eq:piece-decay} to
obtain \eqref{eq:all-pieces-decay}.
\end{proof}

The initial pin mass has now spread sufficiently in mean square.
The remaining term in \eqref{eq:loss-lower} is its pairing with
$g_u$.  We next use flows at the very slice where this pairing occurs,
so that the positive heat energy can absorb it.

\subsection{Absorbing the interaction at the same time}

\begin{proposition}[Expected leakage]\label{prop:leakage}
Fix the constants in the preceding geometric estimates and choose
$C_*$ as in Lemma~\ref{lem:localization}.  For all sufficiently small
$p>0$, every finite cube, every circular grid of mesh at most $p^{4d}$,
every $E$ consisting of $D_0$ and grid pins, and every disjoint
single-time grid test $T$ satisfying \eqref{eq:leakage-growth}, assume
$\mathsf I_p(P)$ for $P=E\cup T$.  If $\beta\ge5s$, then, for
each of the forward and reverse propagations,
\begin{equation}\label{eq:expected-leakage}
 \mathbb E_P\bigl[\|f\|_2^2-\|C_{2s}f\|_2^2\bigr]
             \ge m[1-\varepsilon(p)],
\end{equation}
where one may take
\begin{equation}\label{eq:epsilon-definition}
 \varepsilon(p)=C\left[e^{-cs}
       +R^\alpha\left(s^{-d/2}+\delta+s^{-1}\right)\right],
 \qquad \delta=C R^d p^{2d}+C e^{-cR^\xi}.
\end{equation}
All constants are independent of the volume, period, grid, pin sets,
and final $p$.
\end{proposition}

\begin{proof}
For $s\le u\le2s$, let
\[
 \mu_u=\sum_{i:\,t(u)\text{ regular for }i}\eta_{i,u}.
\]
By \eqref{eq:local-pieces}, $0\le\mu_u\le\eta_u$.
The nonnegative discarded mass is bounded by the localization error
plus the masses attached to nonregular anchors.  Since $0\le g_u\le1$
and each local mass is at most one, Lemma~\ref{lem:localization} and
\eqref{eq:leakage-delta} give
\begin{equation}\label{eq:discarded-pairing}
 \mathbb E_P\langle\eta_u-\mu_u,g_u\rangle
                         \le C m(e^{-cs}+\delta).
\end{equation}
This uses only pointwise bounds, so dependence of $g_u$ on the same
environment causes no change.

At a fixed regular slice, each summand of $\mu_u$ meets exactly the
support and mass hypotheses of \eqref{eq:leakage-flow-input}.  Choose
a reference orientation for the available edges, with
$\operatorname{div}J$ equal to outgoing minus incoming flow.  Since
$g_u$ is finitely supported, summation by parts has finitely many
nonzero terms and gives
\[
 \langle\mu_u,g_u\rangle
   =\sum_{e=(j,k)}J(e)(g_u(j)-g_u(k)).
\]
Cauchy--Schwarz, \eqref{eq:leakage-flow-input}, and
$\sum_e|\nabla_e g_u|^2=2\mathcal D_{t(u)}(g_u)$ imply the
deterministic inequality
\begin{equation}\label{eq:same-time-pairing}
       |\langle\mu_u,g_u\rangle|
          \le\bigl(A\mathcal D_{t(u)}(g_u)\bigr)^{1/2},
               \qquad A=C m R^\alpha.
\end{equation}
It suffices that a flow exists at each time: the resulting inequality
contains only measurable propagated vectors and energies, so no
measurable or continuous choice of flows is required.

To make the remainder in weighted Young's inequality integrable at the
final time, use an averaging density that vanishes linearly there; its
tail weight then vanishes quadratically.
Average \eqref{eq:loss-lower} over $u\in[s,2s]$ with probability
density
\[
 w(u)=\frac{2(2s-u)}{s^2},\qquad
 W(u)=\int_u^{2s}w(v)\,\mathrm dv=\frac{(2s-u)^2}{s^2}.
\]
Since $\mathcal L$ is nondecreasing,
$\mathcal L(2s)\ge\int_s^{2s}w(u)\mathcal L(u)\,\mathrm du$.
Tonelli's theorem shows that the positive energy in this average
contains
\[
                  2\int_s^{2s}W(u)\mathcal D_{t(u)}(g_u)\,\mathrm du;
\]
the additional contribution from $[0,s]$ is nonnegative.
At each $s\le u<2s$, weighted Young's inequality gives
\begin{equation}\label{eq:weighted-young}
 2w(u)\sqrt{A\mathcal D_{t(u)}(g_u)}
 \le2W(u)\mathcal D_{t(u)}(g_u)+\frac{A w(u)^2}{2W(u)}.
\end{equation}
The same-time energy has therefore absorbed the retained pairing.
The remaining integral is finite, because
\begin{equation}\label{eq:weight-integral}
                \int_s^{2s}\frac{w(u)^2}{W(u)}\,\mathrm du=\frac4s.
\end{equation}
The endpoint $u=2s$ has measure zero; the quotient has the constant
value $4/s^2$ on $[s,2s)$.

Taking expectation, using \eqref{eq:discarded-pairing}, and then
Lemma~\ref{lem:regular-time-decay}, we conclude that
\begin{align*}
 \mathbb E_P\mathcal L(2s)
 &\ge m
  -C m\bigl[e^{-cs}+R^\alpha(s^{-d/2}+\delta)\bigr]
  -C m(e^{-cs}+\delta)-C mR^\alpha/s\\
 &\ge m\left[1-C\left(e^{-cs}
                   +R^\alpha(s^{-d/2}+\delta+s^{-1})\right)\right].
\end{align*}
Here $R\ge1$ allows the unweighted $\delta$ error to be included
in $R^\alpha\delta$.  Every step applied to either ordered
propagation, which proves both assertions.
\end{proof}

\subsection{Choosing the constants and closing the finite induction}

The error has enough room to recover the assumed sparsity with the
same parameter $p$.  Indeed, after $C_*$ is fixed, $R$ is comparable
to $s=p^{-a}$ for $p\le1$.  The nonexponential errors $R^\alpha/s$,
$R^\alpha s^{-d/2}$, and $R^{\alpha+d}p^{2d}$ in
\eqref{eq:epsilon-definition} are bounded by constant multiples of
\begin{equation}\label{eq:leakage-exponents}
 p^{a(1-\alpha)},\qquad
 p^{a(d/2-\alpha)},\qquad
 p^{2d-a(d+\alpha)},
\end{equation}
respectively.  Their exponents satisfy
\[
 a(1-\alpha)=\frac{1089}{1000}>1,\qquad
 a(d/2-\alpha)\ge\frac{1639}{1000}>1,\qquad
 2d-a(d+\alpha)\ge\frac{2689}{1000}>1.
\]
The last two lower bounds use $d\ge3$; each exponent increases with
$d$.  Both $e^{-cs}$ and $R^\alpha e^{-cR^\xi}$ are smaller than
every fixed power of $p$.  Thus
\begin{equation}\label{eq:epsilon-small}
                         \varepsilon(p)=o(p)\quad(p\downarrow0).
\end{equation}

For clarity, the constant choices can be made in the following order.
Fix $d,\ell,\alpha,\nu,a$.  Fix the witness scale $b$ and then its
growth constant $L\ge\max\{2d,\ell,4\}$ as in
Lemma~\ref{lem:component-tail}.  Fix the local connectivity radii
and a sufficiently large $M_*$ for
Proposition~\ref{prop:local-nash} and the hull and flow constructions
of Lemma~\ref{lem:hulls} and Proposition~\ref{prop:flow-decay}.
Fix a preliminary $p_0>0$ that absorbs the admissible witness count
in Lemma~\ref{lem:component-tail}; retain the resulting fixed
constants in Proposition~\ref{prop:regularity}, including $c$ and
$\xi$.  Next choose $C_*$ for
Lemma~\ref{lem:localization}.  Finally reduce $p\le p_0$ so that
$R$ exceeds the fixed local radii and the fixed lower bounds needed
for the flow comparisons, $p<1/4$, and
$2\varepsilon(p)\le p$.  The limit
\eqref{eq:epsilon-small} permits all of these final requirements
simultaneously.  No constant in the component tail is recomputed
using this final $p$.

\begin{theorem}[Uniform pin estimate]\label{thm:pin-bootstrap}
For every integer $d\ge3$ and integer slot multiplicity $\ell\ge1$, there
exist $L\ge\max\{2d,\ell,4\}$ and $p_*\in(0,1/4)$ with the
following property.  Let $0<p\le p_*$ and
$\beta\ge5p^{-11/10}$.  In any finite cube $\Lambda$, choose a
finite circular time grid containing zero and of mesh at most
$p^{4d}$.  Let $E$ consist of $D_0$ and any subset of the slot
points on this grid.  Then every nonempty test $T$ at a single grid
time, disjoint from $E$ and satisfying
\eqref{eq:leakage-growth}, obeys
\begin{equation}\label{eq:bootstrap-closed}
                  \mathbb P_E(T\subseteq\mathcal B_E)\le p^{|T|}.
\end{equation}
In particular, with $p=p_*$ and
$\beta_0=5p_*^{-11/10}<\infty$, for every finite
$\beta\ge\beta_0$, every cube, and every slot $j$ at time zero,
\begin{equation}\label{eq:singleton-missing}
                  \mathbb P_{D_0}(j\in\mathcal B_{D_0})\le p_*.
\end{equation}
The constants depend only on $d$ and $\ell$.
\end{theorem}

\begin{proof}
Fix all auxiliary constants in the stated order and choose $p_*$ so
that every $0<p\le p_*$ meets the final requirements.  The set of
interior slot points on a fixed finite grid is finite.  Induct on the
number of these points missing from $E$.  When none is missing there
is no disjoint nonempty test, so $\mathsf I_p(E)$ holds vacuously.

Assume the assertion for pin sets with fewer missing grid points than
$E$, and let $T$ be a test for $E$.  Since $T$ is nonempty and
disjoint from $E$, the set $P=E\cup T$ has strictly fewer missing
points.  The induction hypothesis therefore supplies
$\mathsf I_p(P)$, exactly the hypothesis used in
Proposition~\ref{prop:regularity} and
Proposition~\ref{prop:leakage}.  The two leakage estimates and
Lemma~\ref{lem:two-ends} give
\[
             \mathbb E_P\langle f,(I-U)f\rangle
                         \ge m(1-\varepsilon(p)).
\]
By Lemma~\ref{lem:return-energy}, or
\eqref{eq:leakage-hitting-input}, it follows that
$r\le\varepsilon(p)$.  Proposition~\ref{prop:pin-test} now yields
\[
 \mathbb P_E(T\subseteq\mathcal B_E)
        \le(2r)^m\le(2\varepsilon(p))^m\le p^m.
\]
This proves the induction step for every allowed $T$, and finiteness
of the grid completes the induction.  A singleton interior slot test
at time zero satisfies the growth bound, so taking $E=D_0$ gives
\eqref{eq:singleton-missing}.  A boundary slot belongs to $D_0$,
and its missing-cycle probability is zero.  This covers all slots.
\end{proof}

Theorem~\ref{thm:pin-bootstrap} supplies the uniform missing-cycle
bound required for the magnetization argument.  The grid and the
larger pin sets were devices for a finite induction; the conclusion
\eqref{eq:singleton-missing} concerns the original continuous
boundary pins alone.
